\input{preamble}

\begin{document}

\title{Exact Marginal--Regret Frontiers\\for Signed-Prefix Compression}

\author{\IEEEauthorblockN{Vu Phan An Nguyen}
\IEEEauthorblockA{\textit{Rider University} \\
Lawrenceville, NJ, USA}}

\maketitle

\begin{abstract}
A ranking served to a consumer who reads an unknown prefix must be good at every budget
at once. If, in addition, every item must be retained about equally often at every budget,
what does that cost? We answer exactly for the signed-prefix family, in which the regret at
budget $k$ is the excess of $|\sum_{\ell\le k}g_{\pi_\ell}|$ over the best achievable at
that budget. Four results. (i)~\emph{Exact zero-band prices.} With exactly uniform
retention the optimum is $4/m(m+4)$ for even $m$ and $(m+3)/m(m+4)$ for odd $m$, against
the unconstrained $2/(m+4)(m-1)$: the price is bounded by $2$ in even size and grows
linearly in odd size. (ii)~\emph{The frontier and its breakpoint.} For even $m$ the whole
frontier in the band width $\eta$ is
$\max\{2/(m+4)(m-1),\,4(1-\eta)/m(m+4)\}$; for odd $m$ the price vanishes exactly when
$\eta\ge1$, and for every $\eta$ the odd optimum is at least $(m+3-(m+1)\eta)/m(m+4)$, so
it stays a factor $\Omega(m)$ above the free value at any fixed $\eta<1$. (iii)~\emph{A
polynomial-size program with the same value.} The game over all $m!$ rankings has exactly
the value of a flow over $O(m^{2})$ prefix compositions, every per-item band carried.
(iv)~\emph{Robustness beyond the symmetric family.} For arbitrary per-item weights the
optimum moves by at most $3\|g-g'\|_{1}/d$, and an odd-size perturbation of $\ell_{1}$ size
at most $1/2m$ keeps a price ratio of at least $(m-1)/4$. A separate timing factor appears
if the budget is chosen after the ranking is drawn: the free value becomes exactly $m-1$
times larger. Every theorem is stated for the paper's own objects and is machine-checked in
Lean~4; numerical results are reported as such.
\end{abstract}

\begin{IEEEkeywords}
ranking under constraints, prefix regret, linear programming duality, formal verification
\end{IEEEkeywords}

\section{Introduction}
\label{sec:intro}

\subsection{The question}

A served compression is often a \emph{ranking}: the system returns an order, and a
downstream consumer takes a prefix whose length is not known in advance. Fairness or
exposure requirements then constrain the ranking \emph{distribution}---each item should be
retained about equally often at every budget. We ask what that constraint costs when the
ranking must be good at every prefix length simultaneously.

Fix weights $g=(g_{1},\dots,g_{m})$ and a normaliser $D>0$. For a permutation $\pi$ and a
budget $k$, the per-budget optimum and the regret are
\begin{equation}
b_{k}(g)=\min_{|S|=k}\Bigl|\sum_{i\in S}g_{i}\Bigr|,\quad
r_{k}(\pi;g)=\frac{\bigl|\sum_{\ell\le k}g_{\pi_{\ell}}\bigr|-b_{k}(g)}{D}.
\label{eq:regret}
\end{equation}
For a distribution $P$ over rankings, the constrained optimum is
\begin{equation}
V_{g}(\eta)=\inf_{P}\ \max_{0<k<m}\ \mathbb{E}_{\pi\sim P}\,r_{k}(\pi;g)
\label{eq:value}
\end{equation}
subject to the retained-marginal band
\begin{equation}
\Bigl|\Pr\nolimits_{P}\bigl(i\in\mathrm{prefix}_{k}\bigr)-\tfrac{k}{m}\Bigr|\le\tfrac{\eta}{m}
\quad(\text{all }i,\ k\le m).
\label{eq:band}
\end{equation}
The budget is fixed before the draw: the objective is $\max_{k}\mathbb{E}\,r_{k}$, not
$\mathbb{E}\max_{k}r_{k}$. Section~\ref{sec:expost} shows the two differ by a factor $m-1$
already without a band, so they are never interchanged below.

\subsection{Results}

We work on the family $g_{m}$ of~\eqref{eq:family} below and write
$V_{m}(\eta)=V_{g_{m}}(\eta)$ and $V^{\mathrm{free}}_{m}=2/(m+4)(m-1)$.
$V^{\mathrm{free}}_{m}$ is the unconstrained optimum: the band is vacuous for
$\eta\ge m$, and Theorems~\ref{thm:even} and~\ref{thm:odd} give
$V_{m}(\eta)=V^{\mathrm{free}}_{m}$ there.

\begin{enumerate}
\item[1)] \textbf{Exact zero-band prices} (Theorem~\ref{thm:zero}). $V_{m}(0)$ in closed
form at every $m$. The ratio to $V^{\mathrm{free}}_{m}$ is $2(m-1)/m<2$ for even $m$ and
$(m+3)(m-1)/2m$ for odd $m$. Uniform retention is almost free in even size and costs a
linear factor in odd size.
\item[2)] \textbf{The frontier and its breakpoint} (Theorems~\ref{thm:even},
\ref{thm:odd}, \ref{thm:relaxed}). The whole even frontier; the exact odd breakpoint
$\eta_{*}=1$; and an odd lower bound, linear in $\eta$, that keeps the price $\Omega(m)$ at
any fixed $\eta<1$. On an initial segment of $\eta$ that bound is the value, which we
establish at twelve points by exact rational certificates.
\item[3)] \textbf{Exact flow equivalence} (Theorem~\ref{thm:flow}). The value
of~\eqref{eq:value} over all $m!$ rankings equals the value of a program over $O(m^{2})$
prefix compositions, with every per-item band carried in both directions. This is what
makes the frontier computable at all.
\item[4)] \textbf{Robustness beyond the symmetric family} (Theorem~\ref{thm:robust}). For
arbitrary weights $g'\in\mathbb{R}^{m}$, with baseline and normaliser recomputed, the
optimum is Lipschitz in $\|g_{m}-g'\|_{1}$. For odd $m$ a perturbation of $\ell_{1}$ size at
most $1/2m$ keeps a price ratio of at least $(m-1)/4$.
\end{enumerate}

\subsection{Prior work, and what does not follow from it}

Randomised rankings with constrained position or exposure marginals are a studied object.
Singh and Joachims~\cite{singh2018} introduce exposure constraints and optimise a linear
utility over the Birkhoff polytope. Celis, Straszak and Vishnoi~\cite{celis2018} bound
group representation in every prefix. Singh, Kempe and Joachims~\cite{singh2021} contrast
ex-ante and ex-post guarantees, as Freeman, Shah and Vaish~\cite{freeman2020} do in
allocation. We take that framing as given and do not claim the flow technique: reducing a
permutation polytope to a compositions DAG is standard~\cite{schrijver2003}.

What is new is the \emph{law} that the reduction lets us prove, for an objective these works
do not treat. The signed-prefix regret~\eqref{eq:regret} is an absolute prefix sum net of
the per-budget optimum, taken at every budget. It is not linear in the ranking, and the
parity of $k$ enters through $b_{k}$. Table~\ref{tab:prior} states the difference at the
level of theorems.

\begin{table}[h]
\caption{Nearest sources, and a theorem of this paper that does not follow from each.}
\label{tab:prior}
\centering
\footnotesize
\begin{tabular}{@{}p{0.30\columnwidth}p{0.62\columnwidth}@{}}
\toprule
Nearest source & Theorem here that does not follow \\
\midrule
Exposure-constrained rankings~\cite{singh2018} & The exact signed-prefix value at uniform
marginals, with a $\Theta(m)$ parity split (Thm.~\ref{thm:zero}) \\
Prefix-fair ranking~\cite{celis2018} & The exact breakpoint in the band width at every
$m$, and the whole even frontier (Thms.~\ref{thm:even}, \ref{thm:odd}) \\
Ex-ante vs.\ ex-post~\cite{singh2021,freeman2020} & The exact free ex-post value, a factor
$m-1$ above the ex-ante one (Prop.~\ref{prop:expost}) \\
Flow reductions~\cite{schrijver2003} & Value equality with every per-item band carried
through relabelling (Thm.~\ref{thm:flow}) \\
\bottomrule
\end{tabular}
\end{table}

\subsection{What is proved, and how}

Three evidence classes appear and are never interchanged: \PROVED\ (a theorem, stated at the
paper's objects and checked in Lean~4~\cite{demoura2021} against Mathlib~\cite{mathlib2020});
\CERT\ (an exact rational certificate verified by a script); \NUM\ (a floating-point solve or
search). Each theorem below carries one line naming the Lean declaration that states it. The
proof architecture and every supporting declaration are in the supplement. All modules
compile with no \texttt{sorry} and no \texttt{native\_decide}, and their axioms are within
\{\texttt{propext}, \texttt{Classical.choice}, \texttt{Quot.sound}\}. A compiled theorem
certifies its statement under its hypotheses; the hypotheses of each headline are either
the model assumptions displayed in it or discharged, and the supplement shows each is
satisfiable.

\section{Setting}
\label{sec:setting}

The family studied throughout is
\begin{equation}
g_{m}=\bigl(-3,\underbrace{-1,\dots,-1}_{p},\underbrace{1,\dots,1}_{q},3\bigr),
\quad D=m+4,
\label{eq:family}
\end{equation}
with $p=\lceil(m-2)/2\rceil$ and $q=\lfloor(m-2)/2\rfloor$. The total is $0$ for even $m$
and $-1$ for odd $m$, and $\|g_{m}\|_{1}=m+4=D$. Every weight is an odd integer, so a prefix
sum of length $k$ has the parity of $k$ and cannot be smaller than $k\bmod2$ in absolute
value. For $g_{m}$ the per-budget optimum is exactly that:
\begin{equation}
b_{k}(g_{m})=k\bmod 2\qquad(0<k<m),
\label{eq:baseline}
\end{equation}
attained by an explicit subset at each budget. So~\eqref{eq:regret} on $g_{m}$ is the
normalised parity excess, and the physical objective of Section~\ref{sec:robust}, with the
baseline recomputed from its definition, agrees with it on $g_{m}$.

\leanbox{\lean{familyBaselineIdentification\_holds} in \lean{LeanPaperCBaselineAttainV1};
\lean{trueValue\_family} in \lean{LeanPaperCStabilityFamilyV1}}

The structure of $g_{m}$ is what the parity split lives on: two heavy items of weight
$\pm3$, and light items of weight $\pm1$ whose counts differ by one in odd size. A $+3$ step
between two budgets of zero excess is impossible, because a zero-excess prefix has absolute
value at most $1$. That single observation drives the odd lower bounds.

\section{Exact zero-band prices}
\label{sec:zero}

\begin{theorem}[Zero-band value]
\label{thm:zero}
For every even $m\ge4$ and every odd $m\ge3$,
\begin{equation}
V_{m}(0)=
\begin{cases}
4/m(m+4), & m\text{ even},\\
(m+3)/m(m+4), & m\text{ odd},
\end{cases}
\label{eq:vzero}
\end{equation}
and the infimum is attained by an explicit law on all $m!$ permutations of $g_{m}$. Hence
$V_{m}(0)/V^{\mathrm{free}}_{m}$ is $2(m-1)/m$ for even $m$ and $(m+3)(m-1)/2m$ for odd $m$.
\end{theorem}

\leanbox{\lean{even\_value\_zero}, \lean{odd\_value\_zero}, \lean{odd\_uniform\_price\_ratio}
in \lean{LeanPaperCEndgameV1}}

\emph{Proof idea.} Each side is one budget and one construction.

For even $m$, the lower side is budget $1$. The weights sum to zero, so the prefix of
length $m-1$ is minus the last item, and $r_{1}+r_{m-1}=(2/D)$ times the number of heavy
items at the two ends. Under exact uniformity each end is heavy with probability $2/m$, so
$\mathbb{E}r_{1}+\mathbb{E}r_{m-1}=8/mD$. The upper side mixes the $m-1$ rankings that insert
the heavy pair $(+3,-3)$ at each position of an alternating base with one more ranking,
$(-3,+1,-1,\dots,+1,-1,+3)$, whose excess is exactly $2$ at every budget. The insertion
rankings contribute a single excess spike each, and the spikes cover the budgets once, so the
mixture is flat at $4/mD$. A sign flip at ratio $1/2$ and uniform relabelling within each
weight class make every per-item marginal exactly $k/m$ without moving the objective.

For odd $m$ the lower side is budget $m-1$, where the baseline vanishes and the loss is
decided by the last item: an item of weight $w$ in last position costs $|1+w|/D$. Uniform
retention makes the last item uniform, and averaging $|1+w|$ over $g_{m}$ gives $(m+3)/m$.
The upper side is the uniform law on the $m$ cyclic shifts of the base ranking
$(+3,-3,-1,+1,\dots,-1)$, relabelled within weight classes. Its expected excess at budget
$k$ is a circular window sum: $(m+4)/m-1$ at odd $k$ and $(k+4)/m$ at even $k$. The maximum
is $(m+3)/m$, at $k=m-1$ only. \hfill$\square$

The two cases differ by a factor $\Theta(m)$ and point in opposite directions. In even size
the heavy items can hide at the ends, where their cost is shared between the first and last
budgets. In odd size the imbalance of the light items leaves a $-1$ that the last budget
must pay, and uniformity forces every item, including $+3$, into last position $1/m$ of the
time.

\section{The frontier and its breakpoint}
\label{sec:frontier}

$V_{m}(\eta)$ is non-increasing and convex in $\eta$: the feasible sets are nested, and a
convex combination of plans for two bands is feasible for the combined band while the
objective, a maximum of linear functionals, is convex. We now determine it.

\begin{theorem}[The even frontier]
\label{thm:even}
For every even $m\ge4$ and every $\eta\ge0$,
\[
V_{m}(\eta)=\max\Bigl\{\frac{2}{(m+4)(m-1)},\ \frac{4(1-\eta)}{m(m+4)}\Bigr\},
\]
so the price vanishes exactly at $\eta_{*}=(m-2)/2(m-1)$.
\end{theorem}

\leanbox{\lean{even\_frontier\_all} in \lean{LeanPaperCEndpointLawV1}; shape:
\lean{value\_antitone\_of\_nested}, \lean{value\_convex\_of\_convex\_feasible} in
\lean{LeanPaperCFrontierV1}}

\begin{theorem}[The odd breakpoint]
\label{thm:odd}
For every odd $m\ge3$ and every $\eta\ge0$, $V_{m}(\eta)=V^{\mathrm{free}}_{m}$ if and only
if $\eta\ge1$.
\end{theorem}

\leanbox{\lean{odd\_breakpoint} in \lean{LeanPaperCEndpointLawV1}}

\begin{theorem}[The odd price survives a band]
\label{thm:relaxed}
For every odd $m\ge3$ and every $\eta\ge0$,
\begin{equation}
V_{m}(\eta)\ \ge\ \frac{m+3-(m+1)\eta}{m(m+4)} .
\label{eq:relaxed}
\end{equation}
In particular, for $\eta\le\eta_{0}<1$, $V_{m}(\eta)\ge(1-\eta_{0})/(m+4)=\Omega(1/m)$, against
$V^{\mathrm{free}}_{m}=\Theta(1/m^{2})$.
\end{theorem}

\leanbox{\lean{relaxed\_value\_odd}, \lean{relaxed\_value\_linear} in
\lean{LeanPaperCRelaxedJoinV1}}

\emph{Proof ideas.} The upper sides of Theorems~\ref{thm:even} and~\ref{thm:odd} are one
construction: insert the heavy pair $(+3,-3)$ at a uniformly random position of the
alternating base. Each such ranking has excess exactly $2$ at one budget and $0$ elsewhere,
so the mixture attains $V^{\mathrm{free}}_{m}$ at every budget. Its retained marginals are
read off by counting: $+3$ is in the prefix of length $k$ with probability $k/(m-1)$ and
$-3$ with probability $(k-1)/(m-1)$. For odd $m$ both lie within $1/m$ of $k/m$, which is
$\eta=1$. For even $m$ a sign flip at ratio $1/2$ averages them to $(2k-1)/2(m-1)$, whose
worst deviation, at the two end budgets, is exactly $\eta_{*}/m$. Relabelling uniformly
within weight classes turns these class marginals into the per-item band~\eqref{eq:band}
on a law over all $m!$ permutations. The segment $4(1-\eta)/m(m+4)$ of the even frontier is
the convex combination of this law with the zero-band law of Theorem~\ref{thm:zero}.

The lower sides use two combinatorial facts about every ranking of $g_{m}$. For even $m$,
$r_{1}+r_{m-1}=(2/D)\,\#\{\text{heavy items at the two ends}\}$. The band bounds each end
probability below by $(1-\eta)/m$: the first position is the budget-$1$ prefix, and the
last is the complement of the budget-$(m-1)$ prefix, which is where the band's upper side
enters. For odd $m$, $\sum_{k}(|S_{k}|-b_{k})\ge2+2[\pi_{m}>0]$: the $+3$ step forces
excess $2$ on the two budgets straddling it, and a positive last item forces two more.
Averaged over budgets, every band-feasible law then exceeds the free value by
$(1-\eta)/Dm$. Because the proofs use the free value only as an upper bound, the margin is
uniform over the feasible set and survives the infimum, with no compactness argument and no
window near the breakpoint.

Theorem~\ref{thm:relaxed} uses only budget $m-1$. The last-position law $\pi_{i}$ satisfies
$|\pi_{i}-1/m|\le\eta/m$ by the band at $k=m-1$, and for any centre $a$,
$\sum_{i}\pi_{i}c_{i}\ge\frac1m\sum_{i}c_{i}-\frac{\eta}{m}\sum_{i}|c_{i}-a|$. With
$c_{i}=|1+g_{i}|$ and $a=2$ the two sums are $m+3$ and $m+1$. The centre matters:
$a=0$ would give $m+3$ in place of $m+1$ and a bound that is never tight.
\hfill$\square$

\paragraph{The odd bound is the value on an initial segment}
Theorem~\ref{thm:relaxed} is not merely a bound. For $m\in\{5,7,9\}$ and
$\eta\in\{0,\tfrac14,\tfrac12,\tfrac34\}$ we exhibit a band-feasible flow with rational
path masses whose objective equals the right-hand side of~\eqref{eq:relaxed} exactly, so at
those twelve points $V_{m}(\eta)=(m+3-(m+1)\eta)/m(m+4)$. The solver only locates the
candidate; feasibility and value are verified in exact rational arithmetic. (\CERT)
Bisecting on the solved frontier, the largest $\eta$ at which equality holds agrees with
$1-2/m(m-1)$ to seven digits at $m=5,\dots,15$ (odd). This is a measured law, not a theorem.
(\NUM) The odd curve strictly between the tight segment and the breakpoint is not
determined here.

\section{The flow program has the same value}
\label{sec:flow}

Enumerating rankings is hopeless past $m=11$. A prefix is characterised, for the objective
and for the band alike after symmetrisation, by its \emph{composition}
\begin{equation}
v=(n_{-3},n_{-1},n_{+1},n_{+3})\le(1,p,q,1),
\label{eq:state}
\end{equation}
of which there are $4(p+1)(q+1)=O(m^{2})$. Let $\mathcal{F}_{m}(\eta)$ be the set of unit
flows on the layered DAG of compositions whose edges add one item, with the band imposed on
the class occupancies.

\begin{theorem}[Flow equivalence]
\label{thm:flow}
For every $m\ge2$, every $\eta$, and every functional $\Phi$ of the budget vector
$(\mathbb{E}r_{k})_{k}$, the infimum of $\Phi$ over band-feasible laws on all $m!$
permutations of $g_{m}$ equals its infimum over $\mathcal{F}_{m}(\eta)$. In particular
$V_{m}(\eta)$ is the value of the flow program, and Theorems~\ref{thm:even}
and~\ref{thm:odd} hold for it verbatim.
\end{theorem}

\leanbox{\lean{actual\_banded\_infima} in \lean{LeanPaperCFlowBandJoinV1};
\lean{value\_eq\_flow}, \lean{even\_frontier\_flow}, \lean{odd\_breakpoint\_flow} in
\lean{LeanPaperCFlowObjectiveV1}}

Three maps make it exact. \emph{Ranking to flow:} a law induces the occupancy of each
composition; conservation is the consistency of the marginals of one joint law, and the
whole budget vector transports. \emph{Symmetrisation:} averaging a law over relabellings
inside each weight class leaves every $r_{k}$ unchanged and replaces each item's marginal by
its class average, which keeps the band. \emph{Flow to ranking:} a unit flow carries a
Markov law on whole trajectories whose layer marginals are its occupancies; a trajectory of
positive mass is a sequence of legal one-item steps, hence a type sequence with the right
multiplicities, hence the layer map of an actual permutation; relabelling uniformly within
each class recovers the per-item band. The two programs reach exactly the same set of budget
vectors, not merely the same infimum, and that set is non-empty for every $\eta\ge0$.

The theorem is about the \emph{value} of the four-type program on $g_{m}$. It gives no
algorithm for arbitrary perturbed weights, which the four-type flow does not represent.
Table~\ref{tab:agree} records the program sizes. Solving it reproduces both anchors
of~\eqref{eq:vzero} and $V^{\mathrm{free}}_{m}$ at $m=5,\dots,11$, and agrees with the
permutation LP at $m=5,\dots,9$ for $\eta\in\{0,\frac14,\frac12,1\}$. (\NUM)

\begin{table}[t]
\caption{Flow-program sizes against the permutation formulation. (\NUM)}
\label{tab:agree}
\centering
\footnotesize
\begin{tabular}{@{}rrrl@{}}
\toprule
$m$ & flow states & flow edges & permutations \\
\midrule
$11$ & $120$ & $316$ & $13{,}860$ type sequences \\
$15$ & $224$ & $612$ & intractable \\
$21$ & $440$ & $1236$ & intractable \\
$31$ & $960$ & $2756$ & intractable \\
\bottomrule
\end{tabular}
\end{table}

\section{Robustness beyond the symmetric family}
\label{sec:robust}

The family $g_{m}$ is highly structured, and a reviewer is entitled to ask whether the
factor-$m$ price lives on exact cancellation. Two different questions must be kept apart.
\emph{Relaxing the constraint} is Theorem~\ref{thm:relaxed}. \emph{Perturbing the weights}
is this section. The two are not combined.

For an arbitrary $g'\in\mathbb{R}^{m}$ let $V(g';\eta)$ be the physical optimum:
\eqref{eq:value} over every permutation law with the per-item band~\eqref{eq:band}, the
baseline $b_{k}(g')$ recomputed from its definition, and $D=\|g'\|_{1}$. Perturbations that
break the equal-weight groups are in scope. By~\eqref{eq:baseline},
$V(g_{m};\eta)=V_{m}(\eta)$.

\begin{theorem}[Robustness]
\label{thm:robust}
Let $m\ge3$, $g'\in\mathbb{R}^{m}$ and $0<d\le\min(m+4,\|g'\|_{1})$, and write
$\varepsilon=\|g_{m}-g'\|_{1}$. Then:
\begin{enumerate}
\item[(a)] for every $\eta\ge0$, $|V(g';\eta)-V_{m}(\eta)|\le3\varepsilon/d$; for even
$m\ge4$ this places $V(g';\eta)$ within $3\varepsilon/d$ of the frontier of
Theorem~\ref{thm:even};
\item[(b)] for odd $m\ge3$, whenever $V(g';m)>0$ and $3\varepsilon/d\le V_{m}(0)$,
\[
\frac{V(g';0)}{V(g';m)}\ \ge\ \frac{V_{m}(0)-3\varepsilon/d}{V^{\mathrm{free}}_{m}+3\varepsilon/d};
\]
\item[(c)] for odd $m\ge3$ and $\varepsilon\le1/2m$, with no further condition,
$V(g';0)/V(g';m)\ge(m-1)/4$.
\end{enumerate}
\end{theorem}

\leanbox{(a) \lean{family\_value\_robust}, \lean{even\_frontier\_robust} in
\lean{LeanPaperCStabilityFamilyV1}, from \lean{trueValue\_stable} in
\lean{LeanPaperCStabilityEndgameV1}; (b), (c) \lean{physical\_price\_ratio},
\lean{physical\_price\_linear} in \lean{LeanPaperCRatioJoinV1}}

Here $V(g';m)$ is the unconstrained optimum of $g'$, since the band is vacuous at
$\eta=m$. The constant $3$ in (a) pays for three moving parts: the prefix sum, the
recomputed baseline, and the normaliser. For (c), take $d=m+3$, which is admissible because
$\|g'\|_{1}\ge\|g_{m}\|_{1}-\varepsilon$. Then (a) at $\eta=m$ gives
$V(g';m)\ge V^{\mathrm{free}}_{m}-3\varepsilon/d>0$, and
$4V_{m}(0)-(m-1)V^{\mathrm{free}}_{m}=(2m+12)/m(m+4)\ge3\varepsilon$ closes the arithmetic.

\paragraph{What the floor cannot say}
Part (b) is a lower bound on the price, and at fixed $\varepsilon$ the floor itself is at
most $V_{m}(0)\,d/3\varepsilon\le2/3\varepsilon$, bounded in $m$. So the theorem certifies a
linear price for perturbations of size $O(1/m)$ and certifies nothing linear beyond that.
That is a limit of the method. It does not show that the price collapses: locating a
collapse needs a perturbation proved to destroy the price, and none is proved.

\paragraph{What a search finds}
We searched for the perturbation that hurts the ratio most. The search used random restarts
followed by coordinate descent, over the $m!$ permutation formulation, with baseline and
normaliser recomputed at every step, at $m=5,6,7$ and $\varepsilon\le1.6$. A search returns
\emph{attained} ratios, which are upper bounds on the worst case. It can exhibit a collapse
but cannot rule one out. At $m=7$ the adversarial search keeps $43.5\%$ of the price at
$\varepsilon=0.8$ and $30.1\%$ at $\varepsilon=1.6$, roughly ten times past the scale $1/2m$
of part~(c). A family that shrinks the heavy anchors keeps $57\%$ at $\varepsilon=1.6$. At
$m=5,6$ no family in the grid halved the price, and a uniform shift at $m=6$ \emph{raises}
it to $2.5\times$. Since the search bounds the worst case from above only, the $m=5,6$
outcome says nothing about robustness at those sizes. (\NUM) Between $\varepsilon=1/2m$,
where a linear price is proved, and $\varepsilon\approx0.8$ at $m=7$, where a perturbation
removing more than half of it is exhibited, nothing is determined.

\section{Choosing the budget after the draw}
\label{sec:expost}

If the budget is chosen after the ranking is drawn, the regret is
$\mathbb{E}_{P}[\max_{k}r_{k}]$. The free problem then has an exact answer, and it is
larger by a factor linear in $m$.

\begin{proposition}[The ex-post free value]
\label{prop:expost}
For every $m\ge3$, every ranking of $g_{m}$ has an interior budget with excess at least $2$,
and some ranking has excess at most $2$ at every budget. Hence, over all laws on the $m!$
permutations with no band,
\[
\inf_{P}\ \mathbb{E}_{P}\Bigl[\max_{0<k<m}r_{k}\Bigr]=\frac{2}{m+4}=(m-1)\,V^{\mathrm{free}}_{m},
\]
and for every law the ex-ante objective is at most the ex-post one.
\end{proposition}

\leanbox{\lean{pathValue\_eq}, \lean{path\_over\_free}, \lean{exists\_defect},
\lean{objective\_le\_pathObjective} in \lean{LeanPaperCExPostV1}}

The lower bound is the $+3$ step again. One heavy item is not last. A budget of excess $0$
has a prefix sum of absolute value at most $1$, so a step of size $3$ cannot join two such
budgets. Excess is even because every weight is odd. The upper bound is attained by the
flat ranking for even $m$ and by the base ranking, whose prefix sums take only the values
$0$, $3$ and $-1$, for odd $m$. This is the free case only. The banded ex-post problem is
open, and the band factor of Theorem~\ref{thm:zero} and this timing factor must not be
multiplied into $m^{2}$.

\section{Evidence}
\label{sec:evidence}

Every number is produced by a script that checks itself against a statement proved in Lean
and refuses to write output if a check fails. Table~\ref{tab:evidence} lists the packages.
Figures~\ref{fig:cprice} and~\ref{fig:cstab} are drawn from stored results without
re-solving anything, after each stored number was checked against the proved formula beside
it.

\begin{table}[t]
\caption{Evidence packages. Each script self-checks against a Lean statement.}
\label{tab:evidence}
\centering
\footnotesize
\resizebox{\columnwidth}{!}{%
\begin{tabular}{@{}llr@{}}
\toprule
Package & Script & Checks \\
\midrule
Flow program, breakpoints & \texttt{EXP\_paperC\_flow\_lp} & 5/5 \\
Frontier, stability & \texttt{EXP\_unified\_structural\_frontier} & 9/9 \\
Breakpoint certificates & \texttt{EXP\_paperC\_breakpoint\_certs} & 5/5 \\
Attainment, tight segment & \texttt{EXP\_paperC\_dual\_certificates} & 7/7 \\
Exact Farkas roots & \texttt{EXP\_paperC\_exact\_dual} & 5/5 \\
Breakpoint law, $m=4,\dots,24$ & \texttt{EXP\_paperC\_general\_m} & 9/9 \\
Stored data vs.\ theorems & \texttt{FIG\_paperC\_endpoints} & 10/10 \\
Perturbation & \texttt{EXP\_paperC\_perturbation} & 9/9 \\
Adversarial price search & \texttt{EXP\_paperC\_robustness} & 4/4 \\
Window sums, value & \texttt{EXP\_paperC\_window} & 7/7 \\
\bottomrule
\end{tabular}}
\end{table}

\begin{figure*}[t]
\centering
\includegraphics[width=0.92\textwidth]{paperC_price_boundary.pdf}
\caption{The exact price and the phase boundary. Top left: for even $m$ the whole frontier
$V_{m}(\eta)/V^{\mathrm{free}}_{m}=\max\{1,\,2(1-\eta)(m-1)/m\}$ (Theorem~\ref{thm:even},
lines); the float LP at $m=6,8$ lies on it to $10^{-17}$ (open markers, \NUM). Top right: for
odd $m$ only the zero-band value (stars) and the breakpoint $\eta_{*}=1$ (bar) are theorems;
the dashed curves between them are LP solutions, not a proved frontier. Bottom left: the
zero-band price by parity, $(m+3)(m-1)/2m$ for odd $m$ against $2(m-1)/m$ for even $m$
(Theorem~\ref{thm:zero}). Bottom right: the breakpoint. Lines are the theorems; the flow-LP
bisection points come from a selected set of sizes ($5$--$15$, $17$, $19$, $21$, $25$,
$31$), not every integer through $31$; the diamonds are exact rational Farkas certificates
at $m=5,\dots,11$ (\CERT). Stored numbers were checked against the theorems before drawing
(\texttt{FIG\_paperC\_endpoints}, 10/10).}
\label{fig:cprice}
\end{figure*}

\begin{figure*}[t]
\centering
\includegraphics[width=0.92\textwidth]{paperC_stability.pdf}
\caption{Robustness, from stored perturbations; nothing was re-solved. Left: all 220
perturbations (160 group perturbations over the flow LP, 60 per-item perturbations over the
$m!$ permutation LP) of the physical objective, against the envelope $3\|g_{m}-g'\|_{1}/d$
of Theorem~\ref{thm:robust}(a), $d=\min(\|g_{m}\|_{1},\|g'\|_{1})$. None enters the shaded
region, which would refute the theorem; the worst ratio is $0.32$. The producer skips solver
failures without recording them, so row counts were compared with its fixed grid: none was
dropped. Centre: the same drifts divided by the envelope. Objectives with the normaliser or
the baseline frozen are drawn only to show which object is perturbed. Right: price ratios
\emph{attained} by the adversarial search (markers), upper bounds on the worst case, against
the fixed-normaliser floor $1/(4\varepsilon+4/(m-1))$ (lines, odd $m$; its status is stated
in the supplement). No attained ratio is below the floor, and how far the markers sit above
it measures nothing about the floor's slack.}
\label{fig:cstab}
\end{figure*}

\section{Limitations}
\label{sec:limits}

\begin{enumerate}
\item[1)] \textbf{One family.} Every exact result is for $g_{m}$ of~\eqref{eq:family}.
Theorem~\ref{thm:robust} leaves it only through a Lipschitz envelope.
\item[2)] \textbf{The odd curve between its anchors.} For odd $m$ we prove the zero-band
value, the breakpoint and the linear lower bound~\eqref{eq:relaxed}; the value between the
tight segment and $\eta=1$ is not determined, and the tight segment itself is exact only at
twelve certified points.
\item[3)] \textbf{Value, not algorithm.} Theorem~\ref{thm:flow} equates two values on
$g_{m}$. It gives no polynomial-size program for perturbed weights, and the LP entries
remain floating-point solutions.
\item[4)] \textbf{The collapse scale.} Theorem~\ref{thm:robust} certifies a linear price only
for perturbations of size $O(1/m)$. Where the price actually collapses is measured only as
an upper bound, at three sizes.
\item[5)] \textbf{The ex-post objective.} For $\mathbb{E}\max_{k}r_{k}$ only the free value is
proved. Its price under a band is unknown, and the band and timing factors are not
multiplied.
\end{enumerate}

\section{Conclusion}

Requiring uniform retention of a served ranking at every budget costs a bounded factor in
even size and a linear factor in odd size, and both prices are exact. The even frontier is
known in closed form, the odd price vanishes exactly at $\eta=1$, and before that it stays
linear at any fixed band width below one. The value over all $m!$ rankings is the value of
an $O(m^{2})$ flow, so the frontier can be traced far past enumeration. The price is stable
under arbitrary per-item perturbations of the weights, and it provably survives
perturbations of size $O(1/m)$. What remains open is named rather than hidden: the odd curve
between its anchors, the scale at which the price collapses, and the banded ex-post problem.

\section*{Companion material}
The supplement gives the full proof architecture of each theorem, every supporting Lean
declaration, and the frozen build it was checked in. The claim map records, for each
headline, its model and quantifiers, its declaration, its remaining assumptions, and its
evidence. Corrections made during preparation are recorded in the project ledger, not in
this text.

\begin{thebibliography}{00}

\bibitem{singh2018} A. Singh and T. Joachims, ``Fairness of exposure in rankings,'' in
\emph{Proc. 24th ACM SIGKDD Int. Conf. Knowledge Discovery and Data Mining}, 2018,
pp.~2219--2228.

\bibitem{celis2018} L. E. Celis, D. Straszak, and N. K. Vishnoi, ``Ranking with fairness
constraints,'' in \emph{Proc. 45th Int. Colloq. Automata, Languages, and Programming
(ICALP)}, 2018.

\bibitem{singh2021} A. Singh, D. Kempe, and T. Joachims, ``Fairness in ranking under
uncertainty,'' in \emph{Advances in Neural Information Processing Systems 34}, 2021,
pp.~11896--11908.

\bibitem{freeman2020} R. Freeman, N. Shah, and R. Vaish, ``Best of both worlds: Ex-ante and
ex-post fairness in resource allocation,'' in \emph{Proc. 21st ACM Conf. Economics and
Computation (EC)}, 2020.

\bibitem{schrijver2003} A. Schrijver, \emph{Combinatorial Optimization: Polyhedra and
Efficiency}. Berlin, Germany: Springer, 2003.

\bibitem{mathlib2020} The mathlib Community, ``The Lean mathematical library,'' in
\emph{Proc. 9th ACM SIGPLAN Int. Conf. Certified Programs and Proofs (CPP)}, 2020,
pp.~367--381.

\bibitem{demoura2021} L. de Moura and S. Ullrich, ``The Lean 4 theorem prover and
programming language,'' in \emph{Proc. 28th Int. Conf. Automated Deduction (CADE)}, 2021,
pp.~625--635.

\end{thebibliography}

\end{document}
